\textbf{Executed comparison.} Positive $\Delta Q=Q_{\rm base}-Q_{\rm case}$ denotes a better full-native-data score. Each row below refits the same two global RV/amplitude parameters; the alternative pixels also use their own refitted baseline. Noise hyperparameters are fixed. Values are rounded ppm, and the original fixed-baseline alternative residuals in the introduction therefore differ slightly from this table.

\begin{center}\small\begin{tabular}{lrrr}\toprule

Shape & Primary B/C/R & Alternate B/C/R & $\Delta Q$, primary/alternate\\\midrule

Baseline & +591/-585/-216 & +424/-463/-183 & -0.00/-0.00\\

Upper source $\times$0.5 & +444/-221/-294 & +280/-121/-245 & +23.73/+19.59\\

Inner $\times$1.5, upper $\times$0.5 & +416/-262/-264 & +260/-168/-213 & +27.28/+19.01\\

Upper opacity $\times$2.014, fixed emission & +464/-168/-330 & +277/-63/-273 & +13.15/+11.81\\

Free model 1, original shape & +184/-165/-184 & +36/-82/-117 & +21.83/-0.52\\

Free model 1, refitted shape & +168/-371/-113 & +33/-282/-42 & +37.75/+2.46\\

Fitted smooth longitude map & +578/-592/-200 & +414/-471/-168 & +0.46/-0.97\\

Localized source patch $\times$4 & +476/-606/-15 & +329/-494/+12 & -13.58/-22.49\\

Upper LOS offset +5 km/s & +362/-565/-99 & +240/-473/-49 & +20.40/-7.83\\

Upper radial speed +13 km/s & +864/-623/-162 & +680/-512/-157 & -72.83/-59.08\\

\bottomrule\end{tabular}\end{center}

\textbf{Source and opacity.} Reducing upper emission improves the trough, but appreciable blue residual remains. Doubling upper opacity at fixed source gives only a small response; increasing the lower-state reservoir while keeping emission approximately fixed behaves much like lowering the source. A stronger inner source also adds red emission. Matching only the selected blue/core means with two source coefficients requires large changes and creates substantial collateral errors; it is not an admissible all-data solution.

\textbf{Kinematics and longitude.} A smooth positive longitude map improves primary $Q$ by less than one. A localized bright region can increase blue emission but barely deepens the core and worsens the global comparison. A uniform upper LOS offset is not a radial wind: the actual ray-projected radial acceleration moves the blue/core residual in the wrong direction. The LOS offset's global preference also changes with pixel selection.

\textbf{Stellar/calibration controls.} With the original planetary RV/amplitude held fixed, the full-data stellar-width drift changes the first night's width by at most 0.64\%; its exact target residuals are approximately $(+551,-548,-185)$ ppm and its primary $\Delta Q$ is 22.18. This removes only about 6.5\% of the blue-minus-core discrepancy. Depth pulsation and broad-curvature drift remove less than 1\%. A target-selected 8.45\% width drift can mimic much of the local shape, but its exact full-data $Q$ worsens by about 1471. It is not a viable correction supported by these data.

\textbf{Model 1 is informative, but not a physical identification.} Its original shape reduces both target residuals much more than uniform source scaling. Reoptimizing its shape further improves global $Q$ on primary pixels while leaving a larger local core error. The alternative-pixel global gain is small or absent. Thus a good local shape, a better global score and a physically realizable atmosphere are distinct requirements.
